% Pillar-1 (scaling) standalone introduction + related work
\section{Introduction}
\label{sec:p1-intro}
Reinforcement learning from verifiable rewards has become a standard final stage
in training reasoning-capable language models, and GRPO~\citep{shao2024deepseekmath, guo2025deepseekr1} is now one of its most widely used algorithms. GRPO replaces
the learned value function of PPO~\citep{schulman2017proximal} with a
group-relative baseline: for each prompt it samples a group of $G$ completions,
scores them with a verifier, and centers the rewards within the group to form
advantages. Because it removes the critic, GRPO is cheap to run at scale, and a
natural question follows: \emph{how does the benefit of GRPO post-training change
as the base model grows?}

Pre-training exhibits smooth power-law scaling in parameter count and data
\citep{kaplan2020scalinglaws, hoffmann2022chinchilla}, and it is tempting to expect
an analogous law for RL post-training. Evidence here is thinner and more
contested~\citep{hilton2023rlscaling, rafailov2024rlhfscaling, nimmaturi2025predictive},
in part because RL gains are entangled with the rollout stack, reward parser,
reference policy, and evaluation protocol, so cross-paper comparisons conflate the
algorithm with its plumbing. This paper isolates the scaling question inside a
single, fixed benchmark.

We study GRPO scaling as one pillar of \ours{}, a controlled benchmark of 70+ runs
across seven RL libraries and five model families ($0.6$B--$\sim$671B) on GSM8K
\citep{cobbe2021training}, HumanEval~\citep{chen2021codex}, and synthetic tool-use
tasks. Rather than assert a new law, we report what the data will and will not
support. Our contributions are: (i) a same-benchmark measurement showing the
cross-scale mean training-reward slope is statistically indistinguishable from zero
over the measured 2--3-decade parameter range; (ii) a pre-registered falsification of a clean
three-phase saturation hypothesis; (iii) an identifiability audit showing no
universal saturation exponent is estimable on this data, so the defensible object
is a local, stack-conditioned saturation law with the endpoint ceiling regressed
on $\log_{10} N$; and (iv) a case study of the Nemotron-120B collapse as a
distinct phase rather than trend noise. We read single-seed frontier-API runs as
illustrative case studies, not universal laws. The title's ``limits'' is therefore
substantive: this paper identifies where a parameter-count scaling law is not
estimable from the corpus; it does not claim that GRPO lacks scaling behavior in
better-controlled data. The program's stronger explanatory candidates are usable
contrastive compute and group-size-conditioned signal, tested separately in the
ZVF and group-size companion papers.

\subsection{Related Work}
\label{sec:p1-related}
Neural scaling laws for pre-training are well established~\citep{kaplan2020scalinglaws, hoffmann2022chinchilla, muennighoff2023datascaling}. Scaling behavior of RL and
RLHF post-training has been examined more recently~\citep{hilton2023rlscaling, rafailov2024rlhfscaling, nimmaturi2025predictive}, but typically within a single
stack, leaving open whether reported trends transfer across libraries and model
families. GRPO and its use in reasoning models originate
with~\citet{shao2024deepseekmath} and~\citet{guo2025deepseekr1};
\citet{ahmadian2024backtobasics} argue that much of the benefit of heavy RL
machinery is recoverable by simpler estimators once the stack is controlled. Our
work complements these by fixing the benchmark and asking, conservatively, which
scaling conclusions survive.
